\chapter{Modèles de taux court}\label{ch:rc:short-rate-models}

Chaque nuit, une banque valorise des dizaines de milliers d'obligations remboursables par anticipation, de swaptions bermudéennes
et de portefeuilles de prêts hypothécaires, chacun étant une option permettant d'exercer à l'une des nombreuses dates
portant sur une courbe entière. Un modèle du seul taux de swap ne suffit pas~: l'exercice à
la troisième année dépend de tous les taux à partir de cette même année. La plus ancienne réponse exploitable
modélise un seul nombre, le taux court instantané, et en déduit la courbe entière
à partir de lui~; la version que la plupart des banques utilisent encore en production s'ajuste à la courbe actuelle
de façon exacte, fournit les prix des obligations et des swaptions sous forme fermée, et se cale sur un arbre qui
valorise une swaption bermudéenne en quelques millisecondes. Ses limites sont aussi connues que ses usages~: un
facteur fait bouger tous les taux ensemble, et une seule volatilité ne peut pas ajuster les
swaptions qu'on lui demandera de couvrir. Ce chapitre construit ce modèle, le
\omterm{def:rc:short-rate-models:hw}{modèle de Hull--White}, l'ajuste à la courbe en euros du chapitre~2 et aux
swaptions co-terminales du chapitre~5, vérifie son arbre par rapport à ses formules, et mesure ce que
qu'apporte un second facteur.

\section{Le taux court et les prix affines des obligations}

\begin{definition}[Modèle de taux court]\label{def:rc:short-rate-models:short}
Un \emph{modèle de taux court}\index{modèle de taux court} précise la dynamique du
taux instantané $r_t$ sous la probabilité risque-neutre $\mathbb Q$ (livre~4,
chapitre~5), et valorise chaque obligation zéro-coupon au moyen de~:
$P(t,T) = \E^{\mathbb Q}_t\bigl[\exp(-\int_t^Tr_s\,ds)\bigr]$.
\end{definition}

\begin{definition}[Modèle affine de structure par terme]\label{def:rc:short-rate-models:affine}
Un \omterm{def:rc:short-rate-models:short}{modèle de taux court} est un \emph{modèle affine de structure par terme}\index{modèle affine de structure par terme}
si ses prix d'obligations sont de forme exponentielle-affine en l'état~:
$P(t,T) = \exp\bigl(\mathcal A(t,T)-\mathcal B(t,T)\,x_t\bigr)$ pour un état $x_t$ (le
taux court ou un vecteur de facteurs) et des fonctions déterministes $\mathcal A,
\mathcal B$. D'après la formule de Feynman--Kac (livre~4, chapitre~4), cette propriété tient
chaque fois que la dérive et la variance de l'état sont affines en lui.
\end{definition}

\begin{definition}[Modèle de Vasicek]\label{def:rc:short-rate-models:vasicek}
Le \emph{modèle de Vasicek}\index{modèle de Vasicek} (1977) fait du taux court un
processus d'Ornstein--Uhlenbeck, $dr_t = \kappa(\bar x-r_t)\,dt+\sigma\,dW_t$~: retour
à la moyenne à la vitesse $\kappa$ vers un niveau $\bar x$, incréments normaux, taux négatifs
possibles.
\end{definition}

\begin{proposition}[Prix obligataires de Vasicek]\label{prop:rc:short-rate-models:vasicek}
Avec $B(\tau) = (1-e^{-\kappa\tau})/\kappa$,
\[
P(t,T) = \exp\Bigl(\bigl(\bar x-\tfrac{\sigma^2}{2\kappa^2}\bigr)\bigl(B(\tau)-\tau\bigr)
-\tfrac{\sigma^2}{4\kappa}B(\tau)^2 - B(\tau)\,r_t\Bigr),\qquad \tau = T-t,
\]
et le taux zéro long tend vers $\bar x-\sigma^2/(2\kappa^2)$~: la convexité tire les
rendements longs sous le niveau moyen.
\end{proposition}
\begin{proof}
$\int_t^Tr_s\,ds$ est normale sachant $r_t$, de moyenne $\bar x\tau+(r_t-\bar
x)B(\tau)$ et de variance $\frac{\sigma^2}{\kappa^2}\bigl(\tau-B(\tau)-\frac\kappa2
B(\tau)^2\bigr)$~; en utilisant $\E[e^{-X}] = e^{-\E X+\Var X/2}$.
\end{proof}

Un modèle gaussien à un facteur attribue à chaque taux zéro $z(t,T) = -\ln P/\tau$ la
volatilité normale $\sigma B(\tau)/\tau$. Le retour à la moyenne est donc le levier du modèle
de la \emph{forme} de la courbe de volatilité, et non une prévision sur la direction
des taux. Avec $\kappa=0$, chaque taux zéro a la même volatilité, et avec
$\kappa>0$, les taux longs sont moins volatils que les taux courts
(la \cref{fig:rc:short-rate-models:voltermstructure}).

\begin{omfigure}
\begin{tikzpicture}
\begin{axis}[width=11.5cm,height=5.4cm,
  xlabel={échéance du taux zéro (années)},ylabel={volatilité normale (pb)},
  tick label style={font=\small},label style={font=\small},
  xmin=0,xmax=30,ymin=20,ymax=85,grid=major,grid style={gray!25},
  legend columns=3,legend style={font=\footnotesize,at={(0.5,-0.25)},anchor=north,draw=none,fill=none,/tikz/every even column/.append style={column sep=8pt}},
  table/col sep=comma]
\addplot[omDef,thick] table[x=t,y=k0]{figdata/rates-credit-risk/07-short-rate-models/voltermstructure.csv};
\addplot[omThm,very thick] table[x=t,y=k3]{figdata/rates-credit-risk/07-short-rate-models/voltermstructure.csv};
\addplot[omProp,thick,dashed] table[x=t,y=k10]{figdata/rates-credit-risk/07-short-rate-models/voltermstructure.csv};
\legend{$\kappa=0$,$\kappa=3~\%$,$\kappa=10~\%$}
\end{axis}
\end{tikzpicture}

\omcaption{Volatilité normale des taux zéro, $\sigma B(\tau)/\tau$, dans un modèle gaussien à un facteur
avec $\sigma = 80$ points de base~: le retour à la moyenne amortit la
volatilité des taux longs (69.1 points de base à dix ans avec $\kappa=3~\%$, 50.6
avec $\kappa=10~\%$). Données~: le tutoriel
du chapitre.}\label{fig:rc:short-rate-models:voltermstructure}
\end{omfigure}

\section{Vasicek et Hull--White}

Le \omterm{def:rc:short-rate-models:vasicek}{modèle de Vasicek} comporte trois paramètres et ne peut pas s'ajuster à la courbe actuelle. Hull et
White rendent ce niveau dépendant du temps et le choisissent pour l'ajuster.

\begin{definition}[Modèle de Hull--White]\label{def:rc:short-rate-models:hw}
Le \emph{modèle de Hull--White}\index{modèle de Hull--White} (1990) est
$dr_t = (\vartheta(t)-\kappa r_t)\,dt+\sigma(t)\,dW_t$ avec $\vartheta(t)$ choisi de sorte
que les prix des obligations du modèle égalent les facteurs d'actualisation d'aujourd'hui pour toute échéance.
De manière équivalente, $r_t = f(0,t)+x_t$ avec
\[
dx_t = \bigl(y(t)-\kappa x_t\bigr)dt+\sigma(t)\,dW_t,\quad x_0=0,\qquad
y(t) = \int_0^te^{-2\kappa(t-u)}\sigma(u)^2\,du .
\]
\end{definition}

\begin{proposition}[Obligations et options sur obligations de Hull--White]\label{prop:rc:short-rate-models:hwbond}
Avec $B(t,T) = (1-e^{-\kappa(T-t)})/\kappa$,
\[
P(t,T) = \frac{P(0,T)}{P(0,t)}\exp\Bigl(-B(t,T)\,x_t-\tfrac12B(t,T)^2y(t)\Bigr),
\]
donc $\ln P(T,S)$ suit une loi normale d'écart-type $\sigma_P = B(T,S)\sqrt{y(T)}$, et un
call en $T$ sur l'obligation arrivant à échéance en $S$, de strike $K$, vaut
$P(0,S)\Phi(h)-KP(0,T)\Phi(h-\sigma_P)$ avec $h = \ln\frac{P(0,S)}{KP(0,T)}/\sigma_P+\sigma_P/2$.
\end{proposition}
\admitted\ (un calcul gaussien sous la probabilité forward en $T$~; Brigo et
Mercurio, 2006, chapitre~3.)

\begin{definition}[Modèle de Black--Karasinski]\label{def:rc:short-rate-models:bk}
Le \emph{modèle de Black--Karasinski}\index{modèle de Black--Karasinski} (1991) rend
le logarithme du taux court à retour à la moyenne et normal, $d\ln r_t =
(\vartheta(t)-\kappa\ln r_t)\,dt+\sigma\,dW_t$~: les taux restent positifs, les prix
d'obligations n'ont pas de forme fermée, et le modèle s'utilise sur des arbres. Il a
perdu du terrain lorsque les taux sont devenus négatifs.
\end{definition}

\section{Calibration sur la courbe et sur les swaptions}

La courbe est ajustée par construction. La volatilité doit être ajustée aux options,
et les options qui comptent sont celles dans lesquelles le produit peut être exercé~: pour
une obligation remboursable par anticipation (callable) bermudéenne de dix ans, exerçable chaque année, les swaptions \emph{co-terminales},
à un an sur neuf, à deux ans sur huit, jusqu'à neuf ans sur un.

\begin{definition}[Décomposition de Jamshidian]\label{def:rc:short-rate-models:jamshidian}
La \emph{décomposition de Jamshidian}\index{décomposition de Jamshidian} (1989)
écrit une option sur une obligation à coupons, dans un modèle à un facteur où chaque prix
d'obligation à l'échéance est une fonction décroissante de l'unique état $x_T$, comme un portefeuille
d'options sur ses obligations zéro-coupon~: avec $x^\star$ l'état auquel l'obligation à
coupons vaut le strike, chaque option sur obligation zéro-coupon a pour strike le prix
de cette obligation en $x^\star$.
\end{definition}

Une swaption receveuse de strike $K$ est un call au pair sur une obligation versant $K$
chaque année et un au terme, si bien qu'elle est une somme de calls sur obligations zéro-coupon
évalués par la \cref{prop:rc:short-rate-models:hwbond}.

\omcode{code/firm/shortrate/firm_shortrate.py}{73}{89}{Swaptions européennes dans le modèle de Hull--White par la décomposition de Jamshidian.}\label{lst:rc:short-rate-models:jamshidian}

\begin{method}[Calibration co-terminale]\label{met:rc:short-rate-models:coterminal}
Fixer $\kappa$ (il façonne la structure par terme de la volatilité~; on le règle souvent à partir du
rapport des volatilités de différentes maturités, ou par convention de desk). Pour les
échéances co-terminales $e_1<\dots<e_m$ vers une date finale commune, prendre
$\sigma(t)$ constante par morceaux sur $[e_{i-1},e_i)$ et résoudre chaque niveau à tour de rôle
de sorte que la swaption à la monnaie arrivant à échéance en $e_i$ retrouve sa volatilité
de marché, les niveaux précédents étant maintenus. Vérifier la régularité et le signe de la
$\sigma(t)$ obtenue.
\end{method}

\begin{example}[Un seul sigma ne suffit pas]\label{ex:rc:short-rate-models:fit}
Sur la courbe OIS en euros du chapitre~2 avec $\kappa=3~\%$, les swaptions co-terminales
à dix ans issues du cube du chapitre~5 ont des volatilités normales à la monnaie allant de
60.8 points de base (un an sur neuf) à 91.0 (neuf ans sur un). Le $\sigma$ constant
optimal, 86.0 points de base, donne des volatilités de modèle comprises entre 76.0 et 76.7~: il
manque l'échéance à un an de 15.3 points de base et celle à neuf ans de 14.2
(\cref{fig:rc:short-rate-models:fit}). Une $\sigma(t)$ constante par morceaux reproduit exactement
les neuf, en montant de 68.7 à 116.3 points de base avant de redescendre à 104.2 sur
le dernier intervalle (\cref{fig:rc:short-rate-models:sigma}).
\end{example}

\begin{omfigure}
\begin{tikzpicture}
\begin{axis}[width=11.5cm,height=5.4cm,
  xlabel={échéance de la swaption (années~; le sous-jacent finit à l'année dix)},ylabel={volatilité normale à la monnaie (pb)},
  tick label style={font=\small},label style={font=\small},
  xmin=0.5,xmax=9.5,ymin=55,ymax=95,grid=major,grid style={gray!25},
  legend columns=3,legend style={font=\footnotesize,at={(0.5,-0.25)},anchor=north,draw=none,fill=none,/tikz/every even column/.append style={column sep=8pt}},
  table/col sep=comma]
\addplot[only marks,mark=*,mark size=2.4pt,black] table[x=e,y=market]{figdata/rates-credit-risk/07-short-rate-models/fit.csv};
\addplot[omDef,thick,mark=square*] table[x=e,y=constant]{figdata/rates-credit-risk/07-short-rate-models/fit.csv};
\addplot[omThm,thick,dashed,mark=triangle*] table[x=e,y=piecewise]{figdata/rates-credit-risk/07-short-rate-models/fit.csv};
\legend{marché (cube),Hull--White un $\sigma$,Hull--White $\sigma(t)$}
\end{axis}
\end{tikzpicture}

\omcaption{Swaptions co-terminales à dix ans~: volatilités de marché et ajustements
de Hull--White avec $\kappa=3~\%$. Une seule volatilité donne une droite presque plate au
milieu~; une volatilité constante par morceaux reproduit chaque point. Données~: courbe du
chapitre~2 et cube synthétique du chapitre~5~; le
tutoriel du chapitre.}\label{fig:rc:short-rate-models:fit}
\end{omfigure}

\begin{omfigure}
\begin{tikzpicture}
\begin{axis}[width=11cm,height=5cm,
  xlabel={temps (années)},ylabel={$\sigma(t)$ (pb)},
  tick label style={font=\small},label style={font=\small},
  xmin=0,xmax=10,ymin=60,ymax=120,grid=major,grid style={gray!25},
  table/col sep=comma]
\addplot[omThm,very thick,const plot] table[x=t,y=sigma]{figdata/rates-credit-risk/07-short-rate-models/sigma.csv};
\end{axis}
\end{tikzpicture}

\omcaption{La volatilité de Hull--White constante par morceaux qui reproduit les neuf
swaptions co-terminales. Elle doit monter pour rendre les échéances tardives plus volatiles
que ne le ferait le seul retour à la moyenne, signe que $\kappa=3~\%$ et la structure par terme
du cube ne s'accordent pas. Données~: tutoriel du chapitre.}\label{fig:rc:short-rate-models:sigma}
\end{omfigure}

\section{Arbres}

Une swaption bermudéenne nécessite la valeur de la poursuite à chaque date d'exercice, dans
chaque état~: une induction arrière sur un treillis. L'arbre trinomial du
livre~5, chapitre~22, est construit ici pour l'état du \omterm{def:rc:short-rate-models:hw}{modèle de Hull--White}.

\begin{method}[L'arbre de Hull--White]\label{met:rc:short-rate-models:tree}
(1) Construire un arbre pour $x$ avec $dx = -\kappa x\,dt+\sigma\,dW$~: pas $\Delta t$, écart
$\Delta x = \sigma\sqrt{3\Delta t}$, nœuds $j\Delta x$ pour $|j|\le j_{\max} =
\lceil0.184/(\kappa\Delta t)\rceil$, trois branches issues de chaque nœud avec des
probabilités qui reproduisent l'espérance $-\kappa j\Delta x\Delta t$ et la variance
$\sigma^2\Delta t$, le branchement étant décalé vers l'intérieur en $\pm j_{\max}$. (2) Décaler l'arbre
de $\alpha_i$ à chaque pas afin qu'il reproduise la courbe~: propager vers l'avant les prix
d'Arrow--Debreu $Q_{i,j}$ et poser $\alpha_i = \frac1{\Delta t}\ln\bigl(\sum_j
Q_{i,j}e^{-j\Delta x\Delta t}/P(0,t_{i+1})\bigr)$. Le taux court au nœud $(i,j)$ vaut
$\alpha_i+j\Delta x$.
\end{method}

\omcode{code/firm/shortrate/firm_shortrate.py}{164}{184}{Ajustement de l'arbre sur la courbe par induction avant sur les prix d'Arrow--Debreu.}\label{lst:rc:short-rate-models:tree}

\begin{omfigure}
\begin{tikzpicture}[x=2.3cm,y=0.9cm,font=\small]
\foreach \j in {-1,0,1} {\draw[omDef,thick] (0,0) -- (1,\j);}
\foreach \a in {-1,0,1} {\foreach \j in {-1,0,1} {\draw[omDef,thin] (1,\a) -- (2,\a+\j);}}
\foreach \a in {-2,...,2} {\fill[omThm] (2,\a) circle (2pt);}
\foreach \a in {-1,0,1} {\fill[omThm] (1,\a) circle (2pt);}
\fill[omThm] (0,0) circle (2pt);
\draw[omProp,thick] (2,2) -- (3,2); \draw[omProp,thick] (2,2) -- (3,1); \draw[omProp,thick] (2,2) -- (3,0);
\foreach \a in {0,1,2} {\fill[omProp] (3,\a) circle (2pt);}
\node[anchor=west,font=\footnotesize,align=left] at (3.1,2) {à $j_{\max}$ la ramification\\ se tourne vers l'intérieur};
\node[below=3pt,font=\footnotesize] at (0,-2.2) {$t_0$};
\node[below=3pt,font=\footnotesize] at (1,-2.2) {$t_1$};
\node[below=3pt,font=\footnotesize] at (2,-2.2) {$t_2$};
\node[below=3pt,font=\footnotesize] at (3,-2.2) {$t_3$};
\node[anchor=east,font=\footnotesize] at (-0.1,0) {$x=0$};
\node[anchor=east,font=\footnotesize] at (1.9,2) {$2\Delta x$};
\end{tikzpicture}

\omcaption{Un arbre trinomial de Hull--White~: trois branches issues de chaque nœud, écart
$\Delta x=\sigma\sqrt{3\Delta t}$~; une fois la largeur atteinte $j_{\max}$, les nœuds
extérieurs se ramifient vers l'intérieur, ce qui traduit le retour à la moyenne. Le taux court
en chaque nœud vaut son $x$ augmenté de l'$\alpha_i$ du pas.}\label{fig:rc:short-rate-models:tree}
\end{omfigure}

\begin{example}[Arbre face à la formule]\label{ex:rc:short-rate-models:tree}
Avec $\sigma = 80$ points de base et $\kappa=3~\%$, le payeur cinq-sur-cinq à la monnaie
(forward 2.933~\%, annuité 4.102) vaut 0.026018 par unité de notionnel selon
la formule de Jamshidian. L'arbre à pas mensuels donne 0.026119, puis
à pas de deux semaines 0.026079, à pas hebdomadaires 0.026044 et à pas d'une demi-semaine
0.026016~: l'erreur décroît à peu près proportionnellement au pas. Le facteur d'actualisation à dix ans
est recalculé exactement à chaque taille de pas, par construction.
\end{example}

\section{Deux facteurs, et ce qu'un seul facteur ne peut pas faire}

Dans tout modèle à un facteur, chaque taux est une fonction du même état, de sorte que
les variations de tous les taux sont parfaitement corrélées. Les variations quotidiennes des taux à deux et à dix ans
des rendements du Trésor américain ont présenté une corrélation de 0.769 sur 2016--2026 (chapitre~6)~; un
produit versant sur la différence de deux taux (un CMS spread, chapitre~6) ne vaut rien
pour un modèle à un facteur.

\begin{definition}[Modèle gaussien à deux facteurs]\label{def:rc:short-rate-models:g2}
Le \emph{modèle gaussien à deux facteurs}\index{modèle gaussien à deux facteurs} (G2++) pose
$r_t = \varphi(t)+x_t+y_t$ avec $dx = -\kappa_1x\,dt+\sigma_1dW^1$, $dy =
-\kappa_2y\,dt+\sigma_2dW^2$, $d\langle W^1,W^2\rangle = \rho\,dt$, et $\varphi$ calibrée sur
la courbe. La corrélation des variations des taux zéro de maturités $\tau_1$ et
$\tau_2$ vaut
\[
\frac{\sum_{a,b}b_a(\tau_1)b_b(\tau_2)c_{ab}}{\sqrt{\sum_{a,b}b_a(\tau_1)b_b(\tau_1)c_{ab}\,\sum_{a,b}b_a(\tau_2)b_b(\tau_2)c_{ab}}},
\quad b_a(\tau)=\frac{B(\kappa_a,\tau)}{\tau},
\]
avec $c_{11}=\sigma_1^2$, $c_{22}=\sigma_2^2$, $c_{12}=c_{21}=\rho\sigma_1\sigma_2$.
\end{definition}

\begin{example}[Décorréler la courbe]\label{ex:rc:short-rate-models:g2}
Avec un facteur lent ($\kappa_1=2~\%$, $\sigma_1=75$ points de base), un facteur rapide ($\kappa_2=50~\%$,
$\sigma_2=90$ points de base) et $\rho=-0.8$, la corrélation du modèle entre les taux zéro
à deux et à dix ans vaut 0.774, proche de l'estimation issue du Trésor, et celle entre
deux et trente ans 0.687 (\cref{fig:rc:short-rate-models:g2}). Le $\rho$ fortement
négatif est typique~: le facteur rapide déplace le court terme à l'encontre du facteur
lent pour produire des mouvements de pente.
\end{example}

\begin{omfigure}
\begin{tikzpicture}
\begin{axis}[width=11.5cm,height=5cm,
  xlabel={maturité du second taux zéro (années)},ylabel={corrélation avec le taux à 2 ans},
  tick label style={font=\small},label style={font=\small},
  xmin=0,xmax=30,ymin=0.6,ymax=1.03,grid=major,grid style={gray!25},
  legend columns=2,legend style={font=\footnotesize,at={(0.5,-0.25)},anchor=north,draw=none,fill=none,/tikz/every even column/.append style={column sep=8pt}},
  table/col sep=comma]
\addplot[omThm,very thick] table[x=t,y=g2]{figdata/rates-credit-risk/07-short-rate-models/g2.csv};
\addplot[omDef,thick,dashed] table[x=t,y=one]{figdata/rates-credit-risk/07-short-rate-models/g2.csv};
\legend{modèle gaussien à deux facteurs,tout modèle à un facteur}
\end{axis}
\end{tikzpicture}

\omcaption{Corrélation des variations instantanées du taux zéro à deux ans avec les taux zéro
d'autres maturités~: égale à un dans tout modèle à un facteur, décroissante avec la
maturité dans le modèle à deux facteurs de l'\cref{ex:rc:short-rate-models:g2}. Données~: le
tutoriel du chapitre.}\label{fig:rc:short-rate-models:g2}
\end{omfigure}

\section{Tutoriel~: Hull--White de la courbe à l'arbre}\label{tut:rc:short-rate-models:tutorial}

\begin{tutorial}
\textbf{Objectif.} Calibrer Hull--White sur une courbe et sur des swaptions co-terminales, et
vérifier son arbre par rapport à sa formule fermée. \textbf{État final~:}
\cref{fig:rc:short-rate-models:fit,fig:rc:short-rate-models:sigma} et les nombres de
l'\cref{ex:rc:short-rate-models:tree}.
\begin{tutsteps}
\item \textbf{Le modèle}~: \texttt{HullWhite(curve, kappa, sigmas, knots)} recalcule
chaque facteur d'actualisation de la courbe €STR du chapitre~2.
\item \textbf{Calibrer}~: \texttt{calibrations()} ajuste un $\sigma$ unique et un
$\sigma(t)$ constant par morceaux sur les neuf co-terminales~;
\texttt{fit\_table()} convertit les prix du modèle en volatilités normales.
\item \textbf{Arbre}~: \texttt{tree\_check(dt=\ldots)} pour des pas allant des mensuels aux demi-hebdomadaires.
\item \textbf{Deux facteurs}~: \texttt{g2\_table()}~; \texttt{fig\_rc\_shortrate.py}
produit les graphiques.
\end{tutsteps}
\textbf{Pour aller plus loin.} Calibrer avec $\kappa=0$ et $\kappa=10~\%$ et comparer
les formes de $\sigma(t)$~; ajouter un $\sigma$ dépendant du temps à l'arbre en faisant varier
la taille du pas.
\end{tutorial}

\section{Construction~: le moteur de taux court}\label{bld:rc:short-rate-models:shortrate}

\begin{build}
\bfield{Rôle} Le modèle qui se trouve derrière les produits remboursables par anticipation du livre~: swaptions bermudéennes
(chapitre~9), prêts hypothécaires (chapitre~12), simulation de l'exposition (chapitre~17).
\bfield{Interface} \texttt{HullWhite(curve, kappa, sigmas, knots)} avec
\texttt{bond}, \texttt{bond\_vol}, \texttt{zbo}, \texttt{swaption},
\texttt{annuity\_forward}~; \texttt{calibrate\_coterminal}~; \texttt{HWTree(curve,
kappa, sigma, horizon, dt)} avec \texttt{step\_back}, \texttt{zero\_bonds\_at},
\texttt{european\_swaption}~; \texttt{vasicek\_bond}, \texttt{zero\_rate\_vol},
\texttt{g2\_correlation}.
\bfield{Règles} La courbe est un objet quelconque disposant de \texttt{df\_t}~; les swaps à
accrues annuelles unitaires en temps modèle~; numpy uniquement.
\bfield{Tests d'acceptation} \texttt{code/firm/shortrate/tests/}~: les obligations reproduisent la
courbe et les options sur obligations satisfont la parité~; payeur moins receveur égale
le swap forward~; l'arbre reproduit les facteurs d'actualisation et converge vers le prix de
Jamshidian lorsque le pas diminue~; la calibration par morceaux reproduit ses cibles~; le taux long de Vasicek
vaut $\bar x-\sigma^2/2\kappa^2$~; G2++ avec un facteur désactivé a une
corrélation égale à un.
\bfield{Pour aller plus loin} Un $\sigma$ dépendant du temps sur l'arbre~; une formule de
swaption G2++~; calibrer $\kappa$ sur le rapport entre volatilités co-terminales et volatilités de colonne.
\end{build}

\begin{omsources}
\item O.\ Vasicek, «~An equilibrium characterization of the term structure~»,
\emph{Journal of Financial Economics} 5(2), 1977.
\item J.\ Hull and A.\ White, «~Pricing interest-rate-derivative securities~»,
\emph{Review of Financial Studies} 3(4), 1990.
\item F.\ Jamshidian, «~An exact bond option formula~», \emph{Journal of Finance}
44, 1989.
\item F.\ Black and P.\ Karasinski, «~Bond and option pricing when short rates are
lognormal~», \emph{Financial Analysts Journal} 47(4), 1991.
\item D.\ Brigo and F.\ Mercurio, \emph{Interest Rate Models: Theory and Practice},
deuxième éd., Springer, 2006, chapitres 3--4.
\end{omsources}

\section{Exercices}

\begin{exercise}[$\star$]\label{exo:rc:short-rate-models:1}
Dans un modèle gaussien à un facteur avec $\sigma = 80$ points de base et $\kappa = 3~\%$,
Calculer la volatilité normale du taux zéro à dix ans et du taux zéro à trente ans.
\end{exercise}

\begin{exercise}[$\star$]\label{exo:rc:short-rate-models:2}
Expliquer pourquoi le \omterm{def:rc:short-rate-models:vasicek}{modèle de Vasicek} ne peut pas reproduire la courbe du jour, et ce que
le \omterm{def:rc:short-rate-models:hw}{modèle de Hull--White} modifie pour y remédier.
\end{exercise}

\begin{exercise}[$\star$]\label{exo:rc:short-rate-models:3}
Dans un \omterm{def:rc:short-rate-models:vasicek}{modèle de Vasicek} avec $\bar x = 3\%$, $\sigma = 1\%$ et $\kappa = 10\%$,
Calculer le taux zéro de maturité longue.
\end{exercise}

\begin{exercise}[$\star\star$]\label{exo:rc:short-rate-models:4}
Pour l'arbre à pas mensuels et $\kappa=3~\%$, calculer $j_{\max}$ et le nombre de
nœuds au pas le plus large. Indiquer comment ce nombre évolue avec des pas hebdomadaires.
\end{exercise}

\begin{exercise}[$\star\star$]\label{exo:rc:short-rate-models:5}
Expliquer pourquoi la \omterm{def:rc:short-rate-models:jamshidian}{décomposition de Jamshidian} échoue dans un modèle à deux facteurs.
\end{exercise}

\begin{exercise}[$\star\star$]\label{exo:rc:short-rate-models:6}
Lire la \cref{fig:rc:short-rate-models:sigma}. Indiquer l'allure que prendrait le
$\sigma(t)$ calibré avec $\kappa=10~\%$, et en donner la raison.
\end{exercise}

\begin{exercise}[$\star\star\star$]\label{exo:rc:short-rate-models:7}
\emph{Programmation.} Avec \texttt{tree\_check}, tabuler l'erreur de l'arbre par rapport au prix de
Jamshidian pour des pas mensuels, bihebdomadaires, hebdomadaires et semi-hebdomadaires, puis
estimer l'ordre de convergence.
\end{exercise}

\begin{exercise}[$\star\star\star$]\label{exo:rc:short-rate-models:8}
\emph{Trouver la faille.} «~Notre \omterm{def:rc:short-rate-models:hw}{modèle de Hull--White} à un facteur reproduit chacune des
swaptions co-terminales, donc il évalue correctement nos \omterm{def:rc:convexity-adjustments-and-constant-maturity-products:spread}{options sur spread CMS}.~» Corriger cette affirmation.
\end{exercise}

\section{Problème~: un seul sigma ne suffit pas}

\begin{problem}[{Problème du week-end --- calibration d'un modèle de Hull--White en production}]\label{pb:rc:short-rate-models:1}
Le traitement de nuit d'une banque évalue des swaptions bermudéennes à dix ans en euros, exerçables
chaque année, avec Hull--White, $\kappa=3~\%$, sur la courbe €STR du chapitre~2. Le validateur du modèle
demande pourquoi le front office utilise un $\sigma(t)$ constant par morceaux alors que la documentation indique
«~Hull--White avec une seule volatilité~».

\medskip\noindent\textbf{Partie~I --- Le marché.}
\begin{enumerate}
\item Indiquer quelles swaptions le modèle doit reproduire, et pourquoi celles-là.
\item Donner leur plage de volatilités et sa forme.
\item Dire ce qu'une volatilité co-terminale croissante révèle sur la volatilité forward.
\item Expliquer pourquoi la courbe n'est pas une cible de calibration.
\item Expliquer pourquoi $\kappa$ est fixé plutôt que calibré.
\end{enumerate}

\medskip\noindent\textbf{Partie~II --- Un seul sigma.}
\begin{enumerate}[resume]
\item Donner le meilleur $\sigma$ constant.
\item Donner la plage de volatilités du modèle.
\item Donner les erreurs à un an et à neuf ans.
\item Indiquer quels exercices de la swaption bermudéenne sont surévalués, et lesquels sous-évalués.
\item Expliquer pourquoi aucun $\sigma$ unique ne permet de corriger cela avec $\kappa=3~\%$.
\end{enumerate}

\medskip\noindent\textbf{Partie~III --- Sigma par morceaux.}
\begin{enumerate}[resume]
\item Donner le premier niveau et le plus haut niveau de $\sigma(t)$.
\item Expliquer pourquoi il chute sur le dernier intervalle.
\item Indiquer si le modèle calibré évalue exactement les swaptions co-terminales, et ce qu'il en est des autres swaptions.
\item Dire ce qu'il suppose sur la volatilité au-delà de la neuvième année.
\item Décrire le risque de calibrer neuf nombres chaque jour.
\end{enumerate}

\medskip\noindent\textbf{Partie~IV --- Jugement.}
\begin{enumerate}[resume]
\item Indiquer ce qu'il faudrait écrire dans la documentation du modèle.
\item Indiquer quel test convaincrait le validateur.
\item Dire ce qu'un modèle à deux facteurs apporterait à ce produit, et ce qu'il coûterait.
\item Énoncer le \emph{résultat nommé}~: les erreurs avec un seul $\sigma$ et la plage de la
$\sigma(t)$ calibrée.
\item En une phrase, indiquer ce qu'un \omterm{def:rc:short-rate-models:short}{modèle de taux court} en production doit reproduire.
\end{enumerate}
\end{problem}

\section{Questions d'entretien}

\begin{interviewq}[$\star$ \iqroles{chercheur, banque}]\label{iq:rc:short-rate-models:1}
Que fait le retour à la moyenne dans le \omterm{def:rc:short-rate-models:hw}{modèle de Hull--White}, en des termes qu'un trader
emploierait~?
\end{interviewq}

\begin{interviewq}[$\star\star$ \iqroles{chercheur}]\label{iq:rc:short-rate-models:2}
Dérivez le prix d'une obligation zéro-coupon dans le \omterm{def:rc:short-rate-models:vasicek}{modèle de Vasicek}.
\end{interviewq}

\begin{interviewq}[$\star\star$ \iqroles{chercheur, développeur}]\label{iq:rc:short-rate-models:3}
Comment calibrez-vous un arbre trinomial sur la \omterm{def:rc:multi-curve-and-collateral-discounting:curves}{courbe d'actualisation} du jour~?
\end{interviewq}

\begin{interviewq}[$\star\star$ \iqroles{trader, chercheur}]\label{iq:rc:short-rate-models:4}
Sur quelles swaptions calibrez-vous un modèle à un facteur pour une swaption bermudéenne
à dix ans, et pourquoi pas toute la matrice~?
\end{interviewq}

\begin{interviewq}[$\star\star\star$ \iqroles{chercheur, risque}]\label{iq:rc:short-rate-models:5}
Que ne peut pas évaluer un modèle à un facteur, et comment savez-vous quand cela a de l'importance~?
\end{interviewq}

\begin{interviewq}[$\star\star\star$ \iqroles{développeur}]\label{iq:rc:short-rate-models:6}
Le prix de votre arbre pour une swaption européenne diffère de la forme fermée de 0.5~\%.
Comment en identifiez-vous la cause~?
\end{interviewq}
